
% biography section
% 


\begin{IEEEbiography}
[{\includegraphics[width=1in,height=1.25in,clip,keepaspectratio]{Xueyong_Zhang.jpg}}]{Xueyong Zhang}
received the B.S. degree in applied physics from Nantong University, China, in 2011, the M.E. degree in microelectronics and solid-state electronics from Southeast University, China, in 2014, and the Ph.D. degree from Nanyang Technological University (NTU), Singapore, in 2022. From 2014-2016, he was a design engineer with Silergy Corp., where he was involved in power management IC design. From 2016-2022, he was a research associate and research fellow with NTU. He is currently a staff engineer with Huawei International Pte Ltd., Singapore. His current research interests include low-power analog IC, neuromorphic VLSI, in-memory computing design, and data converters.
\end{IEEEbiography}


\begin{IEEEbiography}[{\includegraphics[width=1in,height=1.25in,clip,keepaspectratio]{ByungKwon.jpg}}]{Byung-Kwon An}
received the B.S. in Department of Electronics Engineering from Kwangwoon University, Korea, Seoul, South Korea, in 2017, He received the M.S degree in circuits for PCM Memories and OTS device from Hanyang University, Korea, Seoul, South Korea, in 2020. He is currently pursuing the Ph.D. degree in design circuits for RRAM Memories circuits with Nanyang Technological University, Singapore. His interests include design for circuit based on non-volatile memory.
\end{IEEEbiography}


\begin{IEEEbiography}[{\includegraphics[width=1in,height=1.25in,clip,keepaspectratio]{ProfKim.jpg}}]{Tony Tae-Hyoung Kim}
(Senior Member, IEEE) received the B.S. and M.S. degrees in electrical engineering from Korea University, Seoul, Korea, in 1999 and 2001, respectively. He received the Ph.D. degree in electrical and computer engineering from the University of Minnesota, Minneapolis, MN, USA in 2009. From 2001 to 2005, he worked for Samsung Electronics where he performed research on the design of high-speed SRAM memories, clock generators, and IO interface circuits. In 2007 ~ 2009 summer, he was with IBM T. J. Watson Research Center and Broadcom Corporation where he performed research on isolated NBTI/PBTI measurement circuits and SRAM Mismatch measurement test structure, and battery-backed memory design, respectively. In November 2009, he joined Nanyang Technological University where he is currently an associate professor. His current research interests include in-memory computing for edge computing, emerging memory circuit design, energy-efficient circuits and systems for IoT and wearable devices, variation and aging tolerant circuits and systems, and circuit techniques for 3-D ICs .
He received Best Paper Award at 2021 ICCE-Asia, Best Demo Award at 2016 IEEE APCCAS, International Low Power Design Contest award at 2016 IEEE/ACM ISLPED, a best paper award at 2014 and 2011 ISOCC, 2008 AMD/CICC Student Scholarship Award, 2008 Departmental Research Fellowship from U. of Minnesota, 2008 IEEE DAC/ISSCC Student Design Contest Award, 2008 Samsung Humantec Thesis Award (Bronze Prize), 2005 ETRI Journal Paper of the Year Award, 2001 Samsung Humantec Thesis Award (Honor Prize), and 1999 Samsung Humantec Thesis Award (Silver Prize). He is an author/co-author of around 190 journal and conference papers and holds 17 US and Korean patents. He serves as an Associate Editor of IEEE Transactions on Very Large Scale Integration (VLSI) Systems, IEEE Access, Frontiers in Electronics, and IEIE Journal of Semiconductor Technology and Science (JSTS). He has also served as a technical committee member of various conferences such as IEEE Asian Solid-State Circuits Conference (A-SSCC), IEEE International Symposium on Circuits and Systems (ISCAS), IEEE/ACM International Symposium on Low Power Electronics and Design (ISLPED), etc. He is the Chair of IEEE CASS VSA-TC and was the Chair of IEEE SSCS Singapore Chapter in 2015~2016. He is an IEEE SSCS distinguished lecturer. He is a senior member of IEEE.
\end{IEEEbiography}

